% ============================================================
%  SCHOLA DOMESTICA — Magister Mathematica
%  Level 1 | Week 12 | Day 4
%  Topic: Subtraction from 6 and Under — Mixed Consolidation
% ============================================================
\documentclass[12pt, letterpaper]{article}

\usepackage[top=1in, bottom=1in, left=1in, right=1in]{geometry}
\usepackage[T1]{fontenc}
\usepackage[utf8]{inputenc}
\usepackage{lmodern}
\usepackage{microtype}
\usepackage{amsmath}
\usepackage{xcolor}
\usepackage{tikz}
\usepackage{array}
\usepackage{enumitem}
\usepackage{multicol}
\usepackage{mdframed}
\usepackage{fancyhdr}

\definecolor{scholared}{RGB}{160,30,30}
\definecolor{scholagold}{RGB}{180,140,40}
\definecolor{scholacream}{RGB}{253,248,235}
\definecolor{scholagreen}{RGB}{50,110,60}
\definecolor{scholablue}{RGB}{40,80,140}

\pagestyle{fancy}
\fancyhf{}
\renewcommand{\headrulewidth}{0.6pt}
\fancyhead[L]{\small\textcolor{scholared}{\textbf{Schola Domestica — Magister Mathematica}}}
\fancyhead[R]{\small\textcolor{scholared}{Level 1 · Week 12 · Day 4}}
\fancyfoot[C]{\small\textcolor{gray}{— \thepage\ —}}

\newcommand{\answerline}{\underline{\hspace{2.5cm}}}
\newcommand{\biganswerline}{\underline{\hspace{4cm}}}
\newcommand{\numbox}{\tikz{\draw[scholablue,thick](0,0)rectangle(1.2cm,0.85cm);}}

\newmdenv[
  backgroundcolor=scholacream,
  linecolor=scholared,
  linewidth=1.5pt,
  roundcorner=6pt,
  innertopmargin=8pt,
  innerbottommargin=8pt,
  innerleftmargin=10pt,
  innerrightmargin=10pt
]{lessonbox}

\newmdenv[
  backgroundcolor={rgb,255:red,235;green,245;blue,235},
  linecolor=scholagreen,
  linewidth=1.2pt,
  roundcorner=4pt,
  innertopmargin=6pt,
  innerbottommargin=6pt,
  innerleftmargin=8pt,
  innerrightmargin=8pt
]{enrichbox}

\begin{document}

\begin{center}
  {\LARGE\textbf{\textcolor{scholared}{Subtraction from Six:}}}\\[4pt]
  {\Large\textcolor{scholagold}{\textit{Putting It All Together}}}\\[6pt]
  {\small Name: \biganswerline \hspace{1cm} Date: \answerline}
\end{center}

\vspace{4pt}
\noindent\textcolor{scholared}{\rule{\linewidth}{1pt}}
\vspace{4pt}

% IMAGE PROMPT (Miyazaki watercolour style):
% A round-faced cat in an apron sorts six brightly coloured fish into two baskets, removing some
% to a smaller bowl. Kitchen setting, warm lamp light, white paper background, no text.

\begin{lessonbox}
\textbf{All subtraction strategies reviewed:}
\begin{itemize}[leftmargin=1.5em, itemsep=2pt]
  \item $-0$: nothing taken, same number remains
  \item $-1$: count back one
  \item $-2$/$-3$: count back from the whole (use number line or fingers)
  \item \textbf{Fact family}: use the related addition to check — $a - b = c$ because $c + b = a$
\end{itemize}
\end{lessonbox}

\bigskip

% ===== PART I: HORIZONTAL SUBTRACTION =====
\noindent\textcolor{scholared}{\large\textbf{Part I — Quick Facts}} \hfill \textit{(Problems 1–8)}

\bigskip

\noindent
\begin{tabular}{@{}p{3.3cm} p{3.3cm} p{3.3cm} p{3.3cm}@{}}
\textbf{1.}\; $6 - 0 =$ \numbox &
\textbf{2.}\; $5 - 2 =$ \numbox &
\textbf{3.}\; $4 - 3 =$ \numbox &
\textbf{4.}\; $6 - 1 =$ \numbox \\[12pt]
\textbf{5.}\; $3 - 3 =$ \numbox &
\textbf{6.}\; $5 - 0 =$ \numbox &
\textbf{7.}\; $6 - 4 =$ \numbox &
\textbf{8.}\; $5 - 3 =$ \numbox
\end{tabular}

\bigskip\bigskip

% ===== PART II: VERTICAL SUBTRACTION =====
\noindent\textcolor{scholared}{\large\textbf{Part II — Vertical Subtraction}} \hfill \textit{(Problems 9–12)}

\bigskip

\noindent
\begin{tabular}{@{}p{2.0cm} p{2.0cm} p{2.0cm} p{2.0cm}@{}}
\textbf{9.}
$\begin{array}{r} 6 \\ -3 \\ \hline \phantom{0} \end{array}$
&
\textbf{10.}
$\begin{array}{r} 5 \\ -1 \\ \hline \phantom{0} \end{array}$
&
\textbf{11.}
$\begin{array}{r} 4 \\ -2 \\ \hline \phantom{0} \end{array}$
&
\textbf{12.}
$\begin{array}{r} 6 \\ -2 \\ \hline \phantom{0} \end{array}$
\end{tabular}

\bigskip\bigskip

% ===== PART III: ADDITION AND SUBTRACTION PAIRS =====
\noindent\textcolor{scholared}{\large\textbf{Part III — Match the Partners}} \hfill \textit{(Problems 13–16)}

\medskip
\noindent\textit{Write the related addition or subtraction sentence.}

\bigskip

\noindent\textbf{13.} If $4 + 2 = 6$, then $6 - 2 =$ \numbox

\medskip
\noindent\textbf{14.} If $3 + 1 = 4$, then $4 - 1 =$ \numbox

\medskip
\noindent\textbf{15.} If $5 - 3 = 2$, then $2 + 3 =$ \numbox

\medskip
\noindent\textbf{16.} If $6 - 4 = 2$, then $2 + 4 =$ \numbox

\bigskip\bigskip

% ===== PART IV: WORD PROBLEMS =====
\noindent\textcolor{scholared}{\large\textbf{Part IV — Story Problems}} \hfill \textit{(Problems 17–20)}

\medskip

\noindent\textbf{17.} Peter has \textbf{6} carrots. He gives \textbf{3} to his sister. How many does he have left?

\medskip\noindent\hspace{2em} $6 - 3 =$ \numbox\ carrots

\bigskip

\noindent\textbf{18.} Five apples are on a tree. Two fall off. How many stay on the tree?

\medskip\noindent\hspace{2em} $5 - 2 =$ \numbox\ apples

\bigskip

\noindent\textbf{19.} A boy had some toy soldiers. He lost \textbf{2}. Now he has \textbf{3}. How many did he start with?

\medskip\noindent\hspace{2em} \numbox $- 2 = 3$\quad He started with \numbox\ soldiers.

\bigskip

\noindent\textbf{20.} Write a subtraction problem with a difference of 0:

\medskip\noindent\hspace{2em} $\numbox - \numbox = 0$

\bigskip

\begin{enrichbox}
\textbf{\textcolor{scholagreen}{Oral Drill — Subtraction from 6}}\\[4pt]
$6{-}1$ \; $5{-}2$ \; $4{-}3$ \; $6{-}4$ \; $5{-}3$ \; $6{-}2$ \; $4{-}0$ \; $6{-}6$ \; $5{-}5$ \; $6{-}3$\\[2pt]
\textit{Aim: no counting needed — instant recall.}
\end{enrichbox}

\vspace{0.3cm}
\noindent\textcolor{gray}{\small\textit{Ray's Primary Arithmetic — all subtraction facts from 6.}}

\end{document}
